\documentclass[12pt]{article}
\usepackage[utf8]{inputenc}
\usepackage[T1]{fontenc}
\usepackage[brazilian]{babel}
\usepackage{geometry}
\usepackage{amsmath}
\usepackage{listings}
\usepackage{xcolor}
\usepackage{fancyhdr}
\usepackage{titlesec}
\usepackage{enumitem}
\usepackage{tcolorbox}
\usepackage{fontawesome5}

\tcbuselibrary{skins,breakable}

\geometry{margin=1in}
\setlength{\headheight}{15pt}
\setlength{\parindent}{0pt}
\setlength{\parskip}{5pt}

% ── Cores ──
\definecolor{primary}{HTML}{1F4E79}
\definecolor{secondary}{HTML}{0B6E99}
\definecolor{accent}{HTML}{2E86C1}
\definecolor{lightbg}{HTML}{EAF2F8}
\definecolor{codebg}{HTML}{F8F9FA}
\definecolor{codeframe}{HTML}{BDC3C7}
\definecolor{fixcolor}{HTML}{E67E22}
\definecolor{fixbg}{HTML}{FEF5E7}
\definecolor{desafiobg}{HTML}{F5EEF8}
\definecolor{desafiocolor}{HTML}{7D3C98}
\definecolor{dicabg}{HTML}{EAFAF1}
\definecolor{dicacolor}{HTML}{1E8449}
\definecolor{exemplobg}{HTML}{F4F6F7}

% ── Caixas personalizadas ──
\newtcolorbox{instrucaobox}{
    colback=lightbg, colframe=primary,
    fonttitle=\bfseries\color{primary},
    title={\faIcon{info-circle} ~Instruções},
    boxrule=1pt, arc=4pt,
    left=8pt, right=8pt, top=6pt, bottom=6pt,
    breakable
}

\newtcolorbox{fixacaobox}[1][]{
    colback=fixbg, colframe=fixcolor,
    fonttitle=\bfseries\color{fixcolor},
    title={\faIcon{pencil-alt} ~Micro Atividade de Fixação},
    boxrule=0.8pt, arc=4pt,
    left=8pt, right=8pt, top=6pt, bottom=6pt,
    breakable, #1
}

\newtcolorbox{exemplobox}[1][]{
    colback=exemplobg, colframe=accent,
    fonttitle=\bfseries\color{accent},
    title={\faIcon{code} ~#1},
    boxrule=0.8pt, arc=4pt,
    left=6pt, right=6pt, top=4pt, bottom=4pt,
    breakable
}

\newtcolorbox{dicabox}{
    colback=dicabg, colframe=dicacolor,
    fonttitle=\bfseries\color{dicacolor},
    title={\faIcon{lightbulb} ~Dica},
    boxrule=0.8pt, arc=4pt,
    left=8pt, right=8pt, top=4pt, bottom=4pt
}

\newtcolorbox{desafiobox}{
    colback=desafiobg, colframe=desafiocolor,
    fonttitle=\bfseries\color{desafiocolor},
    title={\faIcon{trophy} ~Desafio Bônus},
    boxrule=1pt, arc=4pt,
    left=8pt, right=8pt, top=6pt, bottom=6pt,
    breakable
}

% ── Títulos ──
\titleformat{\section}
  {\Large\bfseries\color{primary}}
  {}
  {0pt}
  {\rule{\linewidth}{0.6pt}\\}

\titleformat{\subsection}
  {\large\bfseries\color{secondary}}
  {}
  {0pt}
  {}

\setlist[itemize]{leftmargin=1.4em,itemsep=2pt,topsep=4pt}
\setlist[enumerate]{leftmargin=1.6em,itemsep=2pt,topsep=4pt}

% ── Cabeçalho/Rodapé ──
\pagestyle{fancy}
\fancyhf{}
\lhead{\textbf{\color{primary}Programação em Python}}
\rhead{\small Prof. Msc. Nator Junior C. da Costa}
\cfoot{\color{gray}\thepage}
\renewcommand{\headrulewidth}{0.4pt}

% ── Código ──
\lstset{
    language=Python,
    basicstyle=\ttfamily\small,
    backgroundcolor=\color{codebg},
    frame=single,
    rulecolor=\color{codeframe},
    keywordstyle=\bfseries\color{blue!80!black},
    commentstyle=\itshape\color{gray},
    stringstyle=\color{red!70!black},
    breaklines=true,
    showstringspaces=false,
    tabsize=4,
    xleftmargin=6pt,
    framexleftmargin=6pt,
    numbers=left,
    numberstyle=\tiny\color{gray},
    numbersep=8pt
}

\begin{document}

% ── Cabeçalho da Lista ──
\begin{center}
\colorbox{primary}{\parbox{\dimexpr\linewidth-2\fboxsep\relax}{%
\centering\color{white}
\vspace{6pt}
{\LARGE \textbf{Lista de Atividades -- Aula 04}}\\[4pt]
{\large Estruturas Condicionais em Python}\\[6pt]
}}
\vspace{8pt}

\begin{tabular}{rl}
\textbf{Disciplina:} & Lógica de Programação em Python\\
\textbf{Professor:} & Prof. Msc. Nator Junior Carvalho da Costa\\
\textbf{Semestre:} & 2026.1\\
\end{tabular}
\end{center}

\vspace{6pt}

% ── Instruções ──
\begin{instrucaobox}
Resolva as atividades abaixo em Python. Nesta aula, pratique:
\begin{itemize}
    \item Estruturas condicionais: \texttt{if}, \texttt{elif}, \texttt{else}
    \item Operadores relacionais: \texttt{>}, \texttt{<}, \texttt{>=}, \texttt{<=}, \texttt{==}, \texttt{!=}
    \item Operadores lógicos: \texttt{and}, \texttt{or}, \texttt{not}
\end{itemize}
Leia cada exemplo com atenção, resolva a \textbf{micro atividade de fixação} logo em seguida e depois avance para a atividade principal.
\end{instrucaobox}

%=============================================================
\section*{Parte 1 -- Decisão Simples (\texttt{if / else})}
%=============================================================

\begin{exemplobox}{Exemplo -- if / else}
\begin{lstlisting}
idade = int(input("Digite sua idade: "))

if idade >= 18:
    print("Maior de idade")
else:
    print("Menor de idade")
\end{lstlisting}
\end{exemplobox}

\begin{fixacaobox}
Sem rodar no computador, diga o que será exibido em cada caso:
\begin{enumerate}
    \item \texttt{idade = 20} $\rightarrow$ Saída: \rule{5cm}{0.4pt}
    \item \texttt{idade = 15} $\rightarrow$ Saída: \rule{5cm}{0.4pt}
    \item \texttt{idade = 18} $\rightarrow$ Saída: \rule{5cm}{0.4pt}
\end{enumerate}
\end{fixacaobox}

\subsection*{Atividade 1: Maior ou Menor de Idade}
Escreva um programa que:
\begin{enumerate}
    \item Leia a idade de uma pessoa com \texttt{input()}
    \item Verifique se a pessoa tem 18 anos ou mais
    \item Exiba \texttt{"Maior de idade"} ou \texttt{"Menor de idade"}
\end{enumerate}

\begin{dicabox}
Use \texttt{int(input(...))} para converter a entrada para número inteiro.
\end{dicabox}

\subsection*{Atividade 2: Número Positivo ou Negativo}
Escreva um programa que:
\begin{enumerate}
    \item Leia um número inteiro
    \item Exiba \texttt{"Positivo"} se for maior que zero, ou \texttt{"Negativo"} caso contrário
\end{enumerate}

\begin{fixacaobox}
Complete o código abaixo:
\begin{lstlisting}
numero = int(input("Digite um numero: "))

if ________ :
    print("Positivo")
________ :
    print("Negativo")
\end{lstlisting}
\end{fixacaobox}

%=============================================================
\section*{Parte 2 -- Múltiplas Condições (\texttt{if / elif / else})}
%=============================================================

\begin{exemplobox}{Exemplo -- if / elif / else}
\begin{lstlisting}
nota = float(input("Digite a nota: "))

if nota >= 7:
    print("Aprovado")
elif nota >= 5:
    print("Recuperacao")
else:
    print("Reprovado")
\end{lstlisting}
\end{exemplobox}

\begin{fixacaobox}
Qual será a saída para cada nota abaixo?
\begin{enumerate}
    \item \texttt{nota = 8.5} $\rightarrow$ \rule{5cm}{0.4pt}
    \item \texttt{nota = 6.0} $\rightarrow$ \rule{5cm}{0.4pt}
    \item \texttt{nota = 3.2} $\rightarrow$ \rule{5cm}{0.4pt}
    \item \texttt{nota = 7.0} $\rightarrow$ \rule{5cm}{0.4pt}
\end{enumerate}
\end{fixacaobox}

\subsection*{Atividade 3: Situação do Aluno}
Escreva um programa que:
\begin{enumerate}
    \item Leia a nota final de um aluno (0 a 10)
    \item Exiba:
    \begin{itemize}
        \item \texttt{"Aprovado"} se nota $\geq$ 7
        \item \texttt{"Recuperação"} se nota $\geq$ 5 e $<$ 7
        \item \texttt{"Reprovado"} se nota $<$ 5
    \end{itemize}
\end{enumerate}

\subsection*{Atividade 4: Par, Ímpar ou Zero}
Escreva um programa que:
\begin{enumerate}
    \item Leia um número inteiro
    \item Se for zero, exiba \texttt{"Zero"}
    \item Se for par, exiba \texttt{"Par"}
    \item Se for ímpar, exiba \texttt{"Ímpar"}
\end{enumerate}

\begin{dicabox}
O operador \texttt{\%} (módulo) retorna o resto da divisão. Se \texttt{n \% 2 == 0}, o número é par.
\end{dicabox}

\begin{fixacaobox}
Corrija o código abaixo (há 2 erros):
\begin{lstlisting}
n = int(input("Numero: "))

if n = 0:
    print("Zero")
elif n % 2 = 0:
    print("Par")
else:
    print("Impar")
\end{lstlisting}
\end{fixacaobox}

%=============================================================
\section*{Parte 3 -- Operadores Lógicos (\texttt{and}, \texttt{or}, \texttt{not})}
%=============================================================

\begin{exemplobox}{Exemplo -- Operadores Lógicos}
\begin{lstlisting}
idade = int(input("Idade: "))
cnh = input("Tem CNH? (s/n): ")

if idade >= 18 and cnh == "s":
    print("Pode dirigir")
else:
    print("Nao pode dirigir")
\end{lstlisting}
\end{exemplobox}

\begin{fixacaobox}
Preencha a tabela com \texttt{True} ou \texttt{False}:
\begin{center}
\begin{tabular}{|c|c|c|c|}
\hline
\texttt{idade} & \texttt{cnh} & \texttt{idade >= 18 and cnh == "s"} & Pode dirigir? \\
\hline
20 & \texttt{"s"} & \rule{2cm}{0.4pt} & \rule{2cm}{0.4pt} \\
\hline
16 & \texttt{"s"} & \rule{2cm}{0.4pt} & \rule{2cm}{0.4pt} \\
\hline
25 & \texttt{"n"} & \rule{2cm}{0.4pt} & \rule{2cm}{0.4pt} \\
\hline
15 & \texttt{"n"} & \rule{2cm}{0.4pt} & \rule{2cm}{0.4pt} \\
\hline
\end{tabular}
\end{center}
\end{fixacaobox}

\subsection*{Atividade 5: Login Simples}
Escreva um programa que:
\begin{enumerate}
    \item Leia \texttt{usuario} e \texttt{senha}
    \item Considere acesso válido apenas se:
    \begin{itemize}
        \item usuário igual a \texttt{"admin"}
        \item \textbf{e} senha igual a \texttt{"python123"}
    \end{itemize}
    \item Exiba \texttt{"Acesso permitido"} ou \texttt{"Acesso negado"}
\end{enumerate}

\subsection*{Atividade 6: Pode Dirigir?}
Escreva um programa que:
\begin{enumerate}
    \item Leia a idade e se a pessoa possui CNH (\texttt{s} ou \texttt{n})
    \item Considere que ela pode dirigir somente se:
    \begin{itemize}
        \item idade $\geq$ 18 \textbf{e} resposta da CNH for \texttt{"s"}
    \end{itemize}
    \item Exiba \texttt{"Pode dirigir"} ou \texttt{"Não pode dirigir"}
\end{enumerate}

\begin{fixacaobox}
Qual operador lógico completa cada frase?
\begin{enumerate}
    \item Pode entrar se tem ingresso \rule{2cm}{0.4pt} tem convite. (\texttt{and} / \texttt{or})
    \item Aprovado se nota $\geq$ 7 \rule{2cm}{0.4pt} frequência $\geq$ 75\%. (\texttt{and} / \texttt{or})
    \item Proibido se \rule{2cm}{0.4pt} é maior de idade. (\texttt{not})
\end{enumerate}
\end{fixacaobox}

%=============================================================
\section*{Parte 4 -- Atividades Integradas}
%=============================================================

\subsection*{Atividade 7: Desconto em Compra}
Escreva um programa que:
\begin{enumerate}
    \item Leia o valor total da compra
    \item Aplique as regras:
    \begin{itemize}
        \item 10\% de desconto se valor $>$ 200
        \item 5\% de desconto se valor $>$ 100 e $\leq$ 200
        \item Sem desconto nos demais casos
    \end{itemize}
    \item Exiba o valor original e o valor final com duas casas decimais
\end{enumerate}

\begin{fixacaobox}
Para uma compra de R\$ 150,00, qual desconto será aplicado? \\[4pt]
Valor final = \rule{5cm}{0.4pt}
\end{fixacaobox}

\subsection*{Atividade 8: Classificação por Faixa Etária}
Escreva um programa que:
\begin{enumerate}
    \item Leia a idade de uma pessoa
    \item Classifique em:
    \begin{itemize}
        \item \texttt{"Criança"} se idade $<$ 12
        \item \texttt{"Adolescente"} se idade $\geq$ 12 e $<$ 18
        \item \texttt{"Adulto"} se idade $\geq$ 18 e $<$ 60
        \item \texttt{"Idoso"} se idade $\geq$ 60
    \end{itemize}
\end{enumerate}

%=============================================================
\section*{Desafios}
%=============================================================

\begin{desafiobox}
\subsection*{Desafio 1: Classificação do IMC}
Crie um programa que:
\begin{enumerate}
    \item Leia peso (kg) e altura (m)
    \item Calcule o IMC: $IMC = \dfrac{peso}{altura^2}$
    \item Classifique:
    \begin{itemize}
        \item \texttt{"Abaixo do peso"} se IMC $<$ 18.5
        \item \texttt{"Peso normal"} se 18.5 $\leq$ IMC $<$ 25
        \item \texttt{"Sobrepeso"} se 25 $\leq$ IMC $<$ 30
        \item \texttt{"Obesidade"} se IMC $\geq$ 30
    \end{itemize}
    \item Exiba o IMC com 2 casas decimais e a classificação
\end{enumerate}

\subsection*{Desafio 2: Calculadora Simples}
Crie um programa que:
\begin{enumerate}
    \item Leia dois números e um operador (\texttt{+}, \texttt{-}, \texttt{*}, \texttt{/})
    \item Realize a operação correspondente
    \item Se o operador for \texttt{/} e o segundo número for zero, exiba \texttt{"Erro: divisão por zero"}
    \item Se o operador for inválido, exiba \texttt{"Operador inválido"}
\end{enumerate}
\end{desafiobox}

\end{document}
